\documentclass{situationist-preprint}

\title{Solipsism and the Argument From Disagreement}
\msid{MND-06-a}
\msfamily{MND-06}
\msdiscipline{Philosophy of Mind}
\mstype{Principal manuscript}
\mskeywords{other minds; solipsism; inference to the best explanation; disagreement; argument from analogy; unity of mind}

\begin{document}
\maketitle

\begin{abstract}
The problem of other minds asks what entitles a subject to believe that anyone besides herself has a mind. The argument from analogy, which infers from one's own case, has long been thought too weak, resting as it does on a single instance. We argue that the existence of other minds is established by inference to the best explanation, and that the decisive datum is genuine disagreement: the subject meets assertions that contradict her considered beliefs, that are backed by reasons she had not considered, and that sometimes correct her. We show that solipsism has only two ways of accounting for such episodes. It can treat them as unexplained, in which case it loses to the hypothesis of other minds on explanatory grounds; or it can ascribe the dissenting assertions to a hidden part of the subject's own mind, in which case the part it posits has beliefs, reasons and a point of view contrary to the subject's, and is by the solipsist's own criterion a second mind. Either way solipsism fails, and the inference to other minds stands.
\end{abstract}

\section{The problem}

Each of us is aware of her own thoughts and feelings in a way in which she is not aware of anyone else's. The solipsist concludes that she has no good reason to believe that other minds exist: the bodies she sees and the voices she hears are given to her, but the minds behind them are not. The classical answer is the argument from analogy: I observe that my own behavior is caused by my mental states, I observe similar behavior in other bodies, and I infer that it has similar causes \citep{Mill1865}. \citet{Russell1948} made the inference a postulate of scientific knowledge. The standard objection is that the analogy generalizes from a single case, my own, to millions, and that I can never check the conclusion in any instance; Wittgenstein and \citet{Malcolm1958} added that the argument presupposes that I understand what it would be for another to have a mental state, which on its own premises I cannot. \citet{Hyslop1972} defended the analogy against the first charge; \citet{Pargetter1984} and \citet{Melnyk1994} recast it as an inference to the best explanation, on which the hypothesis that others have minds is accepted because it best explains their behavior. \citet{Avramides2001} surveys the debate.

The explanatory version faces a familiar challenge. Most behavior that the hypothesis of other minds explains can also be described, with some loss of economy, by a hypothesis on which the apparent others are elaborate regularities in the solipsist's experience. The solipsist can grant that the other-minds hypothesis is simpler, and deny that simplicity weighs much against her. We argue that one class of data is not merely explained less economically by solipsism, but cannot be explained by it at all without abandoning it: the data of genuine disagreement. The argument is in the spirit of Davidson's claim that the concept of objective truth is available only to a creature that interacts with another \citep{Davidson1991}, but it does not depend on Davidson's account of interpretation.

The question addressed here is that of the solipsist: whether there are other minds at all. It is distinct from the question how one ordinarily knows what another person feels. PHE-12 argues that empathy gives other minds directly, though not primordially; that is an account of ordinary knowledge of particular minds, and it is compatible with the present result, which concerns the justification available against one who doubts that there is anything for empathy to give.

\section{Framework}

\begin{definition}[Solipsism]\label{d:sol}
\emph{Solipsism} is the thesis that the subject's own mind is the only mind there is. All apparent persons other than the subject are, if anything, contents or regularities of the subject's experience.
\end{definition}

\begin{definition}[Mind]\label{d:mind}
A \emph{mind} is a system of mental states, such as beliefs, experiences and reasonings, that is unified in the sense that its states are jointly available to one point of view: they can be brought together in reasoning, and conflicts among them can be noticed and addressed by the system that has them.
\end{definition}

Definition~\ref{d:mind} records what the solipsist herself means by \enquote{my mind}. Her doubt about other minds is a doubt whether there are points of view other than hers; it would be no comfort to her to learn that there is only one mind in the world if that mind contained two points of view unknown to each other. The criterion of unity is the one defended for consciousness in MND-15, which holds that experience is unified at each time in normal subjects.

\begin{definition}[Genuine disagreement]\label{d:dis}
A subject $S$ is in \emph{genuine disagreement} with a source $A$ about $p$ when (i) $S$ believes $p$ after reflection; (ii) $A$ asserts not-$p$ and maintains the assertion when $S$ presents her reasons; (iii) $A$ supports not-$p$ with reasons that $S$ had not considered and could not have produced from her own beliefs at the time; and (iv) in some such episodes, $S$ comes, on assessing $A$'s reasons, to recognize that not-$p$ is true.
\end{definition}

Clause (iii) excludes the case in which the apparent other merely repeats what $S$ already thinks, and clause (iv) the case in which the dissent is mere noise. Episodes satisfying all four clauses are common: a colleague points out an error in a proof one had checked, a stranger gives directions that contradict one's map and turn out right. They are presupposed by any norm for responding to disagreement, such as the conciliatory norm defended in EPI-07.

\begin{principle}[Inference to the best explanation]\label{p:ibe}
If a hypothesis explains a body of data and its rivals either fail to explain the data or explain them only by positing what the hypothesis posits, one is justified in accepting the hypothesis \citep{Lipton2004}.
\end{principle}

\section{The argument}

\begin{lemma}\label{l:dilemma}
Solipsism can account for genuine disagreement only by (a) treating it as unexplained, or (b) attributing the dissenting assertions to a process within the subject that believes not-$p$, assesses reasons for it, and maintains it against the subject.
\end{lemma}

\begin{proof}
By Definition~\ref{d:sol}, every apparent assertion of not-$p$ is, for the solipsist, an occurrence in $S$'s experience. Either the occurrence has no explanation in terms of anything that represents not-$p$ as true, or it has one. In the first case the episode is unexplained: the fact that the occurrence presents reasons for not-$p$ that $S$ did not have, that persist against her reasons, and that turn out on her own assessment to be correct, is a coincidence. In the second case the explanation must cite something within the only mind there is that represents not-$p$ as true, is sensitive to reasons bearing on it (since it produces reasons new to $S$ and answers hers), and is not moved by $S$'s reasoning (by clause (ii) of Definition~\ref{d:dis}). That is a process that believes not-$p$, assesses reasons and maintains its view against $S$.
\end{proof}

\begin{lemma}\label{l:second}
A process of the kind described in Lemma~\ref{l:dilemma}(b) is a mind distinct from the subject's.
\end{lemma}

\begin{proof}
The process has beliefs and reasonings, and these are unified with one another, since it brings reasons together in support of not-$p$ and answers objections. Its beliefs are not available to $S$'s point of view: $S$ believes $p$ and is unaware of the reasons for not-$p$ until they are presented to her as coming from outside, which is incompatible with their being jointly available to her in reasoning. By Definition~\ref{d:mind}, the process is a mind, and it is not $S$'s.
\end{proof}

\begin{theorem}\label{t:main}
Given genuine disagreement, one is justified in believing that other minds exist, and solipsism cannot explain the existence of genuine disagreement.
\end{theorem}

\begin{proof}
By Lemma~\ref{l:dilemma}, solipsism either leaves genuine disagreement unexplained or explains it by positing a process that, by Lemma~\ref{l:second}, is another mind; in the second case it is no longer solipsism. The hypothesis that there are other minds, some of which disagree with $S$, explains each clause of Definition~\ref{d:dis}: the other believes not-$p$, has reasons of her own, is not moved by $S$'s reasons when they are weaker, and is sometimes right. By Principle~\ref{p:ibe}, $S$ is justified in accepting it.
\end{proof}

The argument differs from the analogical argument in what it takes as evidence. The analogy starts from resemblance between my behavior and another's; the present argument starts from a difference between my beliefs and what another asserts. Resemblance is what solipsism finds easiest to accommodate, since a world of regularities modeled on my own conduct is just what an elaborate dream would contain. A source that knows better than I do is what it cannot accommodate.

\section{Objections and replies}

\begin{objection}[Dreams]\label{o:dream}
In dreams one meets characters who disagree with one and say things one did not expect. Yet no other mind is involved. Disagreement is therefore no evidence of other minds.
\end{objection}

\begin{reply}
Dream interlocutors rarely satisfy clause (iii) and almost never clause (iv): they do not supply reasons the dreamer could not have produced, and the dreamer, on waking, does not find that they were right about anything she had got wrong. Where a dream yields a genuine insight, the episode is the subject's own reasoning, unified with her other states and available to her after waking, and the dream character is its vehicle. Genuine disagreement is a pattern sustained over many episodes, with sources that are reliably better informed than the subject about some matters, and dreams do not exhibit it.
\end{reply}

\begin{objection}[Divided self]\label{o:divided}
The solipsist can say that the hidden process is part of her own mind, which is larger than her conscious point of view. Lemma~\ref{l:second} then begs the question by defining minds as unified.
\end{objection}

\begin{reply}
The solipsist may so extend the word, but the extension gives up what solipsism claimed. Her doubt was whether there is any point of view besides hers; on the extended usage, her \enquote{mind} contains other points of view, with beliefs, reasons and knowledge she lacks. The hypothesis is then a variant of the other-minds hypothesis with a peculiar view of where the other minds are located, and it inherits the burden of explaining why these points of view are reliably associated with distinct bodies, voices and histories.
\end{reply}

\begin{objection}[The demon]\label{o:demon}
A deceiving demon could produce the appearances of disagreement. The argument does not exclude this.
\end{objection}

\begin{reply}
A demon who invents reasons that the subject lacks and maintains views contrary to hers is itself a mind distinct from hers, so the demon hypothesis is not a version of solipsism. Whether it is a rival to the ordinary hypothesis of many embodied minds is a question about external-world scepticism, which EPI-02 answers.
\end{reply}

\section{Conclusion}

Other minds exist, and we are justified in believing so by inference to the best explanation. The decisive data are those of genuine disagreement: assertions contrary to our considered beliefs, backed by reasons we did not have, and sometimes correct. Solipsism must either leave them unexplained or explain them by positing another point of view, which is to posit another mind.

\begin{thebibliography}{99}\small
\bibitem[Avramides(2001)]{Avramides2001} Avramides, A. (2001). \emph{Other Minds}. Routledge.
\bibitem[Davidson(1991)]{Davidson1991} Davidson, D. (1991). Three varieties of knowledge. In A.~P. Griffiths (ed.), \emph{A.~J. Ayer: Memorial Essays}, Royal Institute of Philosophy Supplement 30, 153--166. Cambridge University Press.
\bibitem[Hyslop and Jackson(1972)]{Hyslop1972} Hyslop, A. and Jackson, F. (1972). The analogical inference to other minds. \emph{American Philosophical Quarterly}, 9(2), 168--176.
\bibitem[Lipton(2004)]{Lipton2004} Lipton, P. (2004). \emph{Inference to the Best Explanation}, 2nd edition. Routledge.
\bibitem[Malcolm(1958)]{Malcolm1958} Malcolm, N. (1958). Knowledge of other minds. \emph{The Journal of Philosophy}, 55(23), 969--978.
\bibitem[Melnyk(1994)]{Melnyk1994} Melnyk, A. (1994). Inference to the best explanation and other minds. \emph{Australasian Journal of Philosophy}, 72(4), 482--491.
\bibitem[Mill(1865)]{Mill1865} Mill, J.~S. (1865). \emph{An Examination of Sir William Hamilton's Philosophy}. Longman, Green.
\bibitem[Pargetter(1984)]{Pargetter1984} Pargetter, R. (1984). The scientific inference to other minds. \emph{Australasian Journal of Philosophy}, 62(2), 158--163.
\bibitem[Russell(1948)]{Russell1948} Russell, B. (1948). \emph{Human Knowledge: Its Scope and Limits}. Allen and Unwin.
\end{thebibliography}

\end{document}
